\documentclass[UTF8,fontset=fandol,11pt,a4paper]{ctexart}
\usepackage[a4paper,margin=1.75cm,headheight=15pt]{geometry}
\usepackage{amsmath,amssymb,mathtools}
\usepackage{booktabs,tabularx,array,longtable}
\usepackage{enumitem}
\usepackage[most]{tcolorbox}
\usepackage{tikz}
\usetikzlibrary{arrows.meta,positioning,matrix,calc}
\usepackage{xcolor,fancyhdr,hyperref,microtype,lastpage,graphicx}

\newcolumntype{Y}{>{\raggedright\arraybackslash}X}
\newcolumntype{L}[1]{>{\raggedright\arraybackslash}p{#1}}
\newcolumntype{C}[1]{>{\centering\arraybackslash}p{#1}}
\setlength{\parindent}{2em}
\setlength{\parskip}{0.34em}
\setlength{\emergencystretch}{3em}
\renewcommand{\arraystretch}{1.2}
\setlength{\tabcolsep}{3pt}
\setlist[itemize]{itemsep=0.18em,topsep=0.35em,leftmargin=2em}
\setlist[enumerate]{itemsep=0.18em,topsep=0.35em,leftmargin=2em}

\definecolor{deep}{RGB}{38,58,72}
\definecolor{soft}{RGB}{244,247,248}
\definecolor{greenish}{RGB}{52,91,74}
\definecolor{warnc}{RGB}{136,61,58}
\definecolor{purple}{RGB}{84,72,112}
\definecolor{rulegray}{RGB}{110,116,120}

\hypersetup{colorlinks=true,linkcolor=deep,urlcolor=purple,pdftitle={CO-04 流水线与指令级并行}}
\pagestyle{fancy}
\fancyhf{}
\lhead{\small\color{deep}\textbf{408 · 计算机组成原理 CO-04}}
\rhead{\small\color{deep}Need · Ready · Overlap · Commit}
\cfoot{\small 第 \thepage\ 页 / 共 \pageref{LastPage} 页}
\renewcommand{\headrulewidth}{0.4pt}

\newtcolorbox{corebox}[1]{breakable,colback=soft,colframe=deep,title=\textbf{#1},boxrule=0.8pt,arc=2mm}
\newtcolorbox{methodbox}[1]{breakable,colback=white,colframe=greenish,title=\textbf{#1},boxrule=0.7pt,arc=1.5mm}
\newtcolorbox{warnbox}[1]{breakable,colback=white,colframe=warnc,title=\textbf{#1},boxrule=0.7pt,arc=1.5mm}
\newtcolorbox{examplebox}[1]{breakable,colback=white,colframe=purple,title=\textbf{#1},boxrule=0.7pt,arc=1.5mm}
\newcommand{\kw}[1]{\textbf{\color{deep}#1}}
\newcommand{\code}[1]{\texttt{#1}}

\tikzset{
  co/.style={draw=deep,rounded corners=1.5mm,text width=2.5cm,minimum height=0.9cm,align=center,fill=soft,font=\small,inner sep=3pt},
  state/.style={co,double,double distance=1.2pt,fill=white},
  flow/.style={-{Latex[length=2mm]},thick,draw=deep},
  ref/.style={-{Latex[length=2mm]},dashed,draw=rulegray}
}

\title{\vspace{-1.1cm}\textbf{流水线与指令级并行}\\[0.4em]
\large 多条指令怎样合法重叠而不破坏顺序语义}
\author{408 · 计算机组成原理 Topic CO-04}
\date{}

\begin{document}
\maketitle

\begin{corebox}{母问题：怎样重叠执行，同时保持与 ISA 顺序轨迹等价？}
CO-03 已给出一条指令内部的值依赖和提交边界。CO-04 把多条指令放到同一时间轴上，生成链是
\[
\boxed{
\text{Instruction Dependencies}
\to\text{Stage / Resource Occupancy}
\to\text{Need--Ready Constraints}
\to\text{Legal Overlap}
\to\text{Forward / Stall / Flush}
\to\text{Precise Commit}}
\]
箭头表示约束求解顺序：先承认程序依赖和资源事实，再判断哪些重叠合法；不是先背“停几拍”，再倒推理由。
\end{corebox}

\begin{methodbox}{本册学习范围}
\begin{itemize}
  \item 本册解释怎样把多条指令放进同一时间表，并在资源冲突、数据依赖和分支改变时保持程序顺序语义。
  \item 本册重点解释流水级、延迟、吞吐量、三类冒险、转发、停顿、冲刷、CPI 和精确异常。
  \item 乱序执行、超标量、多线程和多核只作为扩展例子；分析仍从“值何时产生、消费者何时需要、结果何时提交”开始。
\end{itemize}
\end{methodbox}

\section{术语先行：先把时间表中的词说清楚}
\begin{methodbox}{本册的九个关键词}
\begin{itemize}
  \item \kw{流水级：}把一条指令的组合逻辑切开后形成的一段工作，例如取指、译码、执行、访存和写回。
  \item \kw{延迟 latency：}一条指令或一次操作从开始到结果可用所需的时间。
  \item \kw{吞吐量 throughput：}流水线稳定运行后单位时间完成的指令数；它不等于单条指令延迟。
  \item \kw{冒险 hazard：}下一拍本来想推进，但数据、硬件资源或下一条 PC 的条件还不满足。
  \item \kw{RAW/WAR/WAW：}读后写、写后读、写后写三种寄存器依赖；其中 RAW 是“后继要读而前驱还没产生”的真实数据等待。
  \item \kw{转发 forwarding：}把已经产生但尚未写回寄存器堆的值直接送给消费者。
  \item \kw{停顿 stall：}冻结部分流水级，等待资源或数据满足条件。
  \item \kw{冲刷 flush：}清除错误路径或异常路径上的年轻指令，防止它们留下软件可见副作用。
  \item \kw{CPI/ILP：}CPI 是平均每条指令占用的时钟周期数；ILP 是同一指令流中可同时执行的独立指令数量。
\end{itemize}
\end{methodbox}

\tableofcontents
\newpage

\section{流水化从哪里生成}

\subsection{朴素方案与失败}
若一条指令依次经历多个功能阶段，非流水实现会让当前未使用的部件闲置。把长组合路径切成 $k$ 段，并在段间加入流水寄存器后，不同指令可以占用不同阶段。

设各段组合延迟为 $T_i$，流水寄存器开销为 $T_{reg}$，则
\[
T_{clk}\ge \max_i T_i+T_{reg}.
\]
理想情况下 $n$ 条指令总时间为
\[
T_{pipe}=(k+n-1)T_{clk},
\]
因此有限任务的平均完成率为
\[
Throughput_n=\frac{n}{(k+n-1)T_{clk}},
\]
只有当 $n$ 足够大时才趋近稳态峰值 $1/T_{clk}$。这一区分能阻断“流水线吞吐率恒等于主频”的误用：具体执行 100 条指令仍要支付 fill/drain，问稳态上限时才忽略这项有限序列开销。但这整个模型仍只适用于无 hazard、无 miss、每级一拍的理想流水。

\subsection{吞吐率与延迟}
流水线主要提高稳态吞吐率，不保证降低单条指令 latency。更深的流水可能缩短 $T_{clk}$，却增加流水寄存器开销、填充/排空、分支惩罚和旁路复杂度。因此比较性能必须回到
\[
T_{CPU}=IC\times CPI\times T_{clk}.
\]

\section{时间模型：Produced、Ready、Need 与 Commit}

对一个值 $v$，必须至少标四个时刻：
\begin{itemize}
  \item \kw{Produced：}功能部件已经计算出 $v$；
  \item \kw{Ready：}从某条实际路径看，消费者能够获得 $v$；
  \item \kw{Need：}消费者所在阶段必须使用 $v$；
  \item \kw{Commit：}$v$ 或其他效果写入软件可见架构状态。
\end{itemize}

RAW 依赖是否需要停顿取决于
\[
t_{ready}(producer,path)\le t_{need}(consumer).
\]
若不等式失败，必须改变时间（stall）或改变路径（增加 forwarding）；若值尚未 produced，任何旁路都无能为力。

\begin{warnbox}{Ready 不等于 Commit}
forwarding 可以让 ALU 结果在写回前被下一条指令使用，但寄存器堆此时可能仍是旧值。可用时刻回答“谁能拿到”，提交时刻回答“软件可见状态何时改变”。
\end{warnbox}

\section{Stage 与时空图：先声明模型}

经典五级教学模型常写为 IF、ID、EX、MEM、WB，但停顿数依赖以下假设：
\begin{itemize}
  \item 各 stage 的真实职责与输入/输出；
  \item 寄存器堆同周期写/读的先后关系；
  \item 指令存储与数据存储是否共享端口；
  \item 分支在哪一级比较并产生目标；
  \item forwarding 的起点和终点；
  \item memory 是否固定一拍返回。
\end{itemize}

\begin{center}
\small
\begin{tabularx}{\linewidth}{C{1.5cm}YYYYY}
\toprule
周期 & IF & ID & EX & MEM & WB\\
\midrule
1 & $I_1$ & & & &\\
2 & $I_2$ & $I_1$ & & &\\
3 & $I_3$ & $I_2$ & $I_1$ & &\\
4 & $I_4$ & $I_3$ & $I_2$ & $I_1$ &\\
5 & $I_5$ & $I_4$ & $I_3$ & $I_2$ & $I_1$\\
\bottomrule
\end{tabularx}
\end{center}

时空图不是装饰：每个格子同时表示 stage、资源占用和流水寄存器边界。发生 stall 时必须说明哪些前级保持、哪个控制包被清零为 bubble、哪些后级继续推进。

\section{三类 Hazard：三种约束冲突}

\subsection{结构冲突}
同一周期两项工作需要同一不可并用资源，例如 IF 与 MEM 共用单端口存储器。解决路径只有三类：复制/增加端口，以空间换吞吐；错开调度，以周期换资源；或改变组织。是否冲突由题设资源图决定，不能由“无数据依赖”推出无冲突。

\subsection{数据冲突}
\begin{itemize}
  \item \kw{RAW：}消费者读生产者将写的新值，是真依赖；
  \item \kw{WAR：}较年轻指令写得过早，破坏较老指令尚未完成的读；
  \item \kw{WAW：}两个写同一目标的提交次序颠倒。
\end{itemize}

在单发射、按序发射、按序读写的简单五级流水中，通常只有 RAW 形成实际 hazard；WAR/WAW 需要乱序读写、多发射或不同完成延迟才可能出现。这个结论必须绑定模型，不能写成所有流水线的定律。

\subsection{控制冲突}
分支方向/目标尚未确定时，前级已经取入后续指令。等待、静态/动态预测、延迟槽和尽早决策都在交换平均停顿、硬件/编译器复杂度和错误路径成本。无论策略如何，错误路径效果不能提交。

\section{三种控制动作：Forward、Stall、Flush}

\begin{center}
\small
\begin{tabularx}{\linewidth}{L{2.0cm}Y Y Y}
\toprule
动作 & 保留/推进什么 & 阻止什么 & 使用条件\\
\midrule
Forward & 正常推进并改用旁路值 & 消费旧值 & 值已经 produced，且实际路径可达 need 端\\
Stall & 保持 PC/前级状态或局部队列 & 未准备好的消费者前进 & need 早于所有 ready 路径，或资源冲突\\
Flush & 保留正确较老状态 & 错路/异常后的年轻效果提交 & 分支误判、异常、重定向\\
\bottomrule
\end{tabularx}
\end{center}

forwarding 不是“打开一个总旁路”。ALU $\to$ EX、load data $\to$ EX、result $\to$ branch compare、store data $\to$ MEM 是不同路径；题目没给路径，就不能假设存在。

\section{贯穿母例：LOAD 后立即 ADD}

\[
I_1:\ \code{LOAD R1,0(R2)},\qquad
I_2:\ \code{ADD R3,R1,R4}.
\]

在经典五级模型中：$I_1$ 的 EA 在 EX 产生，数据通常在 MEM 末尾产生；$I_2$ 在自己的 EX 开始需要 $R1$。若相邻发射，数据的 ready 晚于 need，必须至少让消费者后移，直到存在可达旁路。

\begin{center}
\small
\begin{tabularx}{\linewidth}{C{1.1cm}YYYYYYY}
\toprule
 & 1 & 2 & 3 & 4 & 5 & 6 & 7\\
\midrule
$I_1$ & IF & ID & EX & MEM & WB & &\\
$I_2$ & & IF & ID & stall & EX & MEM & WB\\
\bottomrule
\end{tabularx}
\end{center}

该“一拍”只在上述 stage、memory latency、寄存器时序和 MEM-to-EX 旁路假设下成立。若分级更深、memory 多周期或旁路端点不同，必须重新比较 need/ready。

\begin{center}
\resizebox{0.96\linewidth}{!}{\providecolor{themebg}{HTML}{FAFAF7}
\providecolor{themefg}{HTML}{111111}
\providecolor{themegray}{HTML}{666666}
\providecolor{themecurve}{HTML}{1D63B8}
\providecolor{themealert}{HTML}{C53030}
\providecolor{themeamber}{HTML}{B86E00}
\begin{tikzpicture}[font=\scriptsize,color=themefg,text=themefg,>=Stealth]
\tikzset{
  cell/.style={draw=themegray,minimum width=9mm,minimum height=6.8mm,inner sep=0pt,align=center},
  active/.style={cell,draw=themecurve,fill=themecurve!10!themebg},
  wait/.style={cell,draw=themealert,fill=themealert!10!themebg},
  bubble/.style={cell,draw=themeamber,dashed,text=themeamber},
  note/.style={rounded corners=3pt,draw=themegray,fill=themefg!3!themebg,inner xsep=6pt,inner ysep=5pt,align=left},
  flow/.style={->,line width=1.05pt,themecurve},
  bad/.style={->,line width=1.05pt,themealert}
}
\node[anchor=west,font=\bfseries\large] at (0,9.35) {经典五段流水线：冒险只在“需要时刻早于可用时刻”时变成停顿};
\node[anchor=west,text=themegray] at (0,8.82)
  {本图固定：IF--ID--EX--MEM--WB；ALU 结果 EX 末产生，Load 数据 MEM 末产生；支持 EX/MEM、MEM/WB $\to$ EX 旁路。};

% time header
\node[anchor=east,font=\bfseries] at (1.45,7.75) {周期};
\foreach \c in {1,...,7}{\node[font=\bfseries,text=themegray] at (1.75+0.92*\c,7.75) {\c};}

% ALU forwarding panel
\node[anchor=west,font=\bfseries,text=themecurve] at (0,7.05) {A. ALU $\to$ ALU：值已产生，旁路即可，不必等 WB};
\node[anchor=east] at (1.45,6.35) {$I_1$: ADD};
\foreach \c/\s in {1/IF,2/ID,3/EX,4/MEM,5/WB}{\node[active] (a1\c) at (1.75+0.92*\c,6.35) {\s};}
\node[anchor=east] at (1.45,5.55) {$I_2$: SUB};
\foreach \c/\s in {2/IF,3/ID,4/EX,5/MEM,6/WB}{\node[active] (a2\c) at (1.75+0.92*\c,5.55) {\s};}
\draw[flow,bend left=24] (a13.north east) to node[above,font=\scriptsize,text=themecurve] {R1 ready} (a24.north west);
\node[note,anchor=west,text width=5.8cm] at (9.3,6.0) {\textbf{Need--Ready：} $I_1$ 的 R1 在周期 3 的 EX 末已产生；$I_2$ 到周期 4 的 EX 才需要，因此旁路路径赶得上。};

% load use
\node[anchor=west,font=\bfseries,text=themealert] at (0,4.65) {B. Load-use：未来值还没产生，旁路也不能“穿越时间”};
\node[anchor=east] at (1.45,3.95) {$I_1$: LW};
\foreach \c/\s in {1/IF,2/ID,3/EX,4/MEM,5/WB}{\node[active] (l1\c) at (1.75+0.92*\c,3.95) {\s};}
\node[anchor=east] at (1.45,3.15) {$I_2$: ADD};
\node[active] (l22) at (1.75+0.92*2,3.15) {IF};
\node[wait] (l23) at (1.75+0.92*3,3.15) {ID};
\node[wait] (l24) at (1.75+0.92*4,3.15) {ID};
\node[active] (l25) at (1.75+0.92*5,3.15) {EX};
\node[active] at (1.75+0.92*6,3.15) {MEM};
\node[active] at (1.75+0.92*7,3.15) {WB};
\node[bubble] at (1.75+0.92*4,2.35) {bubble};
\draw[bad,bend left=21] (l14.south east) to node[below,font=\scriptsize,text=themealert] {data ready at MEM end} (l25.north west);
\node[note,anchor=west,text width=5.8cm] at (9.3,3.55) {原本 $I_2$ 想在周期 4 进入 EX，但 Load 数据到周期 4 \textbf{末}才存在；把 consumer 的 EX 推到周期 5 后，MEM/WB $\to$ EX 才能送达。};

% structural hazard strip
\node[anchor=west,font=\bfseries,text=themeamber] at (0,1.55) {C. 结构冒险是“同拍抢同一资源”，不是数据依赖};
\node[cell,draw=themeamber] (mem) at (3.0,.65) {$I_1$ MEM\\data memory};
\node[cell,draw=themeamber] (fetch) at (6.0,.65) {$I_4$ IF\\instruction fetch};
\node[font=\bfseries,text=themealert] at (4.5,.65) {$\times$};
\node[note,anchor=west,text width=7.3cm] at (8.1,.65) {只有题设让 IF 与 MEM 共用\textbf{同一个单端口存储资源}时才冲突；若 I/D 分离或有足够端口，这条结构冒险不存在。};

\node[anchor=west,font=\scriptsize,text=themegray] at (0,-.45)
  {禁止误读：本图的一拍 load-use stall 只属于上面的经典五段与旁路时序；换 RF 时序、MEM 延迟、旁路端点或乱序模型后必须重新比较 Produced / Path-Ready / Need。};
\end{tikzpicture}
}
\end{center}
图把 ALU 旁路、Load-use 与结构冒险放在同一 Need--Ready/Resource 模型下；一拍 stall 和 IF/MEM 冲突都明确绑定图上写出的经典五段假设。

\begin{warnbox}{母例故意暴露的错误路径}
“有 forwarding，所以永不停顿”忽略了值尚未产生；“LOAD 在 EX 算出了结果”把 EA 当成 load data；“I1 已经转发，所以已经完成”把 ready 当 commit。
\end{warnbox}

\section{控制冒险与预测状态机}

\subsection{惩罚从决策位置生成}
若分支在第 $d$ 级末尾才确定，错误路径已经进入的年轻指令数量由前端组织和预测时机决定。penalty 不是“分支固定几拍”，必须从决策级、已推进阶段和 flush 范围数出。

平均 CPI 贡献可写为
\[
CPI_{branch}=f_{branch}\,r_{mispredict}\,P_{mispredict},
\]
其中 $P_{mispredict}$ 必须是额外周期；若题目给的是总分支周期，需先剥离理想基线，避免重复计数。

\subsection{两位饱和计数器}
两位预测器是四状态 FSM，单次相反结果只使状态向另一方向移动一步。题目必须给初态、实际 T/NT 序列和采用高位/状态预测规则；循环“通常只错一次”不是无条件定律，初态和连续多轮会改变结果。

\begin{center}
\resizebox{0.96\linewidth}{!}{\providecolor{themebg}{HTML}{FAFAF7}
\providecolor{themefg}{HTML}{111111}
\providecolor{themegray}{HTML}{666666}
\providecolor{themecurve}{HTML}{1D63B8}
\providecolor{themealert}{HTML}{C53030}
\providecolor{themeamber}{HTML}{B86E00}
\begin{tikzpicture}[font=\scriptsize,color=themefg,text=themefg,>=Stealth]
\tikzset{
  st/.style={rounded corners=3pt,draw=themegray,fill=themefg!3!themebg,minimum width=2.25cm,minimum height=.85cm,align=center},
  taken/.style={st,draw=themecurve,fill=themecurve!10!themebg},
  nottaken/.style={st,draw=themeamber,fill=themeamber!10!themebg},
  arr/.style={->,line width=.9pt,themegray},
  wrong/.style={rounded corners=2pt,draw=themealert,fill=themealert!10!themebg,minimum width=1.35cm,minimum height=.65cm,align=center},
  ok/.style={rounded corners=2pt,draw=themecurve,fill=themecurve!10!themebg,minimum width=1.35cm,minimum height=.65cm,align=center},
  note/.style={rounded corners=3pt,draw=themegray,fill=themefg!3!themebg,inner xsep=6pt,inner ysep=5pt,align=left}
}
\node[anchor=west,font=\bfseries\large] at (0,8.35) {2-bit 分支预测：状态机决定“猜哪边”，流水线决定“猜错要作废谁”};
\node[anchor=west,text=themegray] at (0,7.82) {计数器只保存近期方向偏好；预测失败的代价由分支实际在哪一级决策、已有多少年轻指令进入流水线决定。};

% FSM
\node[anchor=west,font=\bfseries,text=themecurve] at (0,6.95) {A. 四状态饱和计数器};
\node[nottaken] (sn) at (2.0,5.65) {00\\强不跳};
\node[nottaken] (wn) at (5.1,5.65) {01\\弱不跳};
\node[taken] (wt) at (8.2,5.65) {10\\弱跳转};
\node[taken] (st) at (11.3,5.65) {11\\强跳转};
\draw[arr] (sn) -- node[above]{T} (wn);
\draw[arr,bend left=18] (wn) to node[above]{T} (wt);
\draw[arr] (wt) -- node[above]{T} (st);
\draw[arr,bend left=18] (st) to node[below]{NT} (wt);
\draw[arr,bend left=18] (wt) to node[below]{NT} (wn);
\draw[arr,bend left=18] (wn) to node[below]{NT} (sn);
\draw[arr,loop left] (sn) to node[left]{NT} (sn);
\draw[arr,loop right] (st) to node[right]{T} (st);
\node[note,anchor=west,text width=4.5cm] at (13.1,5.65) {00/01 预测 NT，10/11 预测 T；一次反常结果只把强状态推到弱状态，不立即翻转方向。初始状态必须由题设给定。};

% pipeline flush
\node[anchor=west,font=\bfseries,text=themealert] at (0,3.95) {B. 预测失败：年轻的错误路径指令必须 Flush，不能 Commit};
\node[anchor=west] at (0,3.15) {假设分支 $B$ 在 EX 末才确定实际方向：};
\node[ok] (bif) at (4.0,3.15) {$B$: IF};
\node[ok] (bid) at (5.6,3.15) {ID};
\node[ok] (bex) at (7.2,3.15) {EX};
\draw[arr] (bif)--(bid); \draw[arr] (bid)--(bex);
\node[wrong] (i1) at (5.6,2.15) {错路 $I_1$\\ID};
\node[wrong] (i2) at (4.0,1.25) {错路 $I_2$\\IF};
\draw[->,line width=1.15pt,themealert] (bex.south) |- node[pos=.25,right,font=\scriptsize,text=themealert]{mispredict} (i1.east);
\draw[->,line width=1.15pt,themealert] (bex.south) |- (i2.east);
\node[font=\bfseries,text=themealert] at (8.7,2.0) {Flush $I_1,I_2$};
\node[ok] (target) at (11.0,2.0) {正确目标\\重新 IF};
\draw[->,line width=1.1pt,themecurve] (8.9,2.0)--(target);
\node[note,anchor=west,text width=4.65cm] at (13.1,2.0) {Flush 与 Stall 不同：Stall 是正确路径暂时不能推进；Flush 是已经进入的年轻指令被证明属于错误路径，必须取消其架构副作用。};

\node[anchor=west,font=\scriptsize,text=themegray] at (0,-.05) {禁止误读：图中的“两条错路指令”只对应“EX 末决策 + 单发射经典前端”这一示例。若分支在 ID/MEM 决策、采用 delay slot 或不同取指宽度，惩罚必须重新由时间线计算。};
\end{tikzpicture}
}
\end{center}
状态机只决定预测方向；mispredict penalty 仍必须由实际分支决策级和已经进入流水线的年轻指令数生成。

\section{CPI 与总时间：从时空图计数}

实际 CPI 可以分解为
\[
CPI_{actual}=CPI_{ideal}
+CPI_{struct}+CPI_{data}+CPI_{control}+CPI_{memory}+\cdots.
\]
分项只有在互斥或题设明确时才能直接相加；若一个 memory miss 同时阻止了后续依赖，不能把同一停顿重复归入两项。

有限指令序列的总周期应从时空图第一拍数到最后一条的规定完成/提交点。填充与排空已经包含在图中，不应再机械追加 $k-1$。

\section{精确异常与提交顺序}

精确状态要求：异常指令之前的较老指令效果已经成立；异常指令和更年轻指令的禁止副作用没有提交。顺序流水可以冻结/清零写使能并 flush；更复杂实现把 produced/finished 与 ordered commit 分开。

\[
\boxed{\text{execution may overlap}\quad\land\quad
\text{architectural trace remains recoverable}}
\]

这条不变量连接 CO-03 的单指令 Write Set、CO-04 的多指令年龄顺序和 X-B01 的异常入口。OS handler 不在本册展开。

\section{代价、边界与 ILP 扩展}

\begin{center}
\small
\begin{tabularx}{\linewidth}{L{2.2cm}Y Y Y}
\toprule
机制 & 换来的能力 & 新成本 & 最小失败条件\\
\midrule
更深流水 & 更短候选时钟、更多重叠 & 寄存器开销、分支/依赖周期数增加 & 最慢段/寄存器开销不降\\
Forwarding & 减少已产生值的 RAW stall & 比较器、MUX、布线和时序压力 & 值尚未产生或端点不可达\\
分支预测 & 降低平均控制停顿 & 预测状态、能耗、错误 flush & 低可预测性或高 penalty\\
多发射 & 峰值 IPC 大于 1 & 端口、依赖检查、带宽 & 指令不独立或资源不足\\
VLIW & 编译期显式并行、简化部分硬件 & 代码密度、调度和兼容性压力 & 动态延迟/并行不足\\
\bottomrule
\end{tabularx}
\end{center}

超标量、乱序、寄存器重命名、ROB、VLIW 和多核属于“同一依赖/资源/提交框架怎样扩张”的 Extension。408 主干只保留能解释 WAR/WAW、动态调度和 ordered commit 的最小接口，不把现代实现参数写成固定规律。

\section{Flynn 与多处理器：并行对象先分清}
流水线和超标量主要是在一条指令流内重叠；多处理器则增加执行上下文或指令流。按指令流/数据流分类：
\begin{center}\small
\begin{tabularx}{\linewidth}{L{2.0cm}L{2.2cm}YY}
\toprule
类别 & 指令流/数据流 & 典型结构 & 题面判据\\\midrule
SISD & 单/单 & 单核顺序处理器 & 一条指令处理一个数据序列\\
SIMD & 单/多 & 向量机、向量指令 & 同一控制对多个数据元素执行\\
MISD & 多/单 & 少见的专用容错结构 & 多条指令作用于同一数据流\\
MIMD & 多/多 & 多核、SMP、通用多处理器 & 不同核心可执行不同线程/程序\\
\bottomrule
\end{tabularx}\end{center}
向量处理属于 SIMD 的数据级并行；多核/SMP 属于 MIMD 的共享或分布式执行资源，不能把“多发射每拍多条”直接说成“多核”。硬件多线程则是在一个核心内保存多个线程上下文，在某线程等待或资源空闲时切换发射；它提高资源利用率，不自动增加物理 ALU 数量。

多处理器题仍沿本册三条约束检查：每个执行流的依赖与 ready/need、共享 Cache/总线/内存的资源冲突、最终架构提交是否符合各自 ISA 顺序。题目若进一步问 Cache 一致性、锁或线程调度，应转到相应专题；本节只确定并行层次与硬件接口。

\begin{methodbox}{并行层级的停止条件}
先回答“增加的是数据元素、同一指令流中的发射宽度、线程上下文，还是物理核心/指令流”；再检查共享资源与提交边界。若无法指出新增并行对象，就不能仅凭“同时”判定属于 SIMD、超标量或多核。
\end{methodbox}

\section{五列概念边界}

\begin{center}
\small
\begin{tabularx}{\linewidth}{L{1.8cm}C{0.35cm}L{1.8cm}Y Y}
\toprule
概念 A & $\neq$ & 概念 B & 真正区别与题目信号 & 混淆后果\\
\midrule
Latency & $\neq$ & Throughput & 单条穿越时间 vs 单位时间完成数 & 误称流水必降单条延迟\\
依赖 & $\neq$ & Hazard & 程序语义关系 vs 当前时间/资源下的冲突 & 把所有 RAW 都算固定停顿\\
Produced & $\neq$ & Ready & 已算出 vs 通过实际路径可被消费者拿到 & 假设不存在的旁路\\
Ready & $\neq$ & Commit & 消费者可用 vs 架构状态已更新 & 精确异常和寄存器值判断错误\\
Stall & $\neq$ & Flush & 等待后继续同一路径 vs 作废错误/年轻状态 & 保留了不应提交的控制包\\
ILP & $\neq$ & 多核并行 & 单指令流内部重叠 vs 多核心/任务并行 & 把 OS 调度和 Cache 一致性塞入流水线题\\
\bottomrule
\end{tabularx}
\end{center}

\section{做题控制协议}

\begin{enumerate}
  \item 写 stage 定义、每级资源、寄存器读写时序、memory latency 和已有 forwarding；
  \item 逐条列源/目的寄存器和控制转移，标 producer、consumer；
  \item 对每个相关值标 produced、ready、need、commit；
  \item 先画无冲突时空图，再叠加结构、数据和控制约束；
  \item forwarding 必须写起点、终点和到达时刻；剩余 ready--need 差才用 stall；
  \item 分支写预测状态、决策级和 flush 范围；
  \item 从最终时空图统计总周期/CPI，不重复添加填充或同一停顿；
  \item 用年龄顺序和 Write Set 检查错误/异常路径没有提交。
\end{enumerate}

\begin{methodbox}{压缩成两问}
流水线题先问：\kw{这个值沿现有路径什么时候 ready，消费者什么时候 need？这个副作用在哪个年龄/时钟边界才允许 commit？}
\end{methodbox}

\section{本册与单指令路径的接口}
\begin{corebox}{从单条路径到多条时间表}
本册给每个值标出 produced（产生）、ready（沿实际路径可取得）、consumer need（消费者需要）和 architectural commit（写入软件可见状态）的时刻，再据此判断合法重叠、转发、停顿、冲刷和提交顺序。阶段名称不能替代这些时刻。
\end{corebox}
吞吐量是稳定填满后的单位时间完成量，latency 是单条指令从进入到结果/提交的时间；两者和 clock rate、CPI、分支 penalty、miss stall 一起才形成 Cost。流水线加深或并行度增加不自动降低完整 CPU time。

\section{考纲映射与一页压缩}

\begin{center}
\small
\begin{tabularx}{\linewidth}{L{3.2cm}Y Y}
\toprule
传统考点 & 本册机制位置 & 下游接口\\
\midrule
流水线性能 & stage 延迟、填充/排空、CPU time & Rules 性能工具箱\\
结构/数据/控制冒险 & 资源、need/ready、Next PC & CO-03/CO-B01\\
转发、停顿、冲刷 & 三种控制动作与接口条件 & 具体时空图题\\
分支预测 & 预测 FSM、决策级、mispredict penalty & CO-I01\\
超标量/乱序/VLIW & ILP Extension 与按序提交边界 & 不扩张为多核 Topic\\
\bottomrule
\end{tabularx}
\end{center}

\begin{corebox}{一句话}
流水线把多条指令放到同一时间轴上：先用 need/ready 和资源约束判断哪些重叠合法，再用 forward、stall、flush 修复冲突，最终只允许符合年龄顺序和 ISA 的效果提交。
\end{corebox}

自测：为什么 forwarding 不能消灭所有 RAW？为什么 WAR/WAW 在简单五级顺序流水中通常不形成 hazard？为什么更深流水不保证更快？为什么分支 penalty 必须从决策级和 flush 范围数出？

\section{边界说明：哪些结论必须依赖实现条件}

经典五级流水、某个寄存器堆的半拍读写规则、现代处理器的级数/发射宽度和“所有 RISC/CISC 性能相当”等说法都必须带实现条件。Flynn 分类、多核、乱序和编译器后端只用于说明并行层次，不替代本册的 need/ready/commit 计算。

\section{补充细节：把流水线看成有约束的时间表}

流水线基础、性能、冒险、异常/中断和多处理器问题，都可以统一回接为
\[
\text{Instruction Stream}
\to\text{Stage Resources}
\to\text{Produced / Ready / Need}
\to\text{Forward / Stall / Flush}
\to\text{Ordered Commit}.
\]
流水线不是“把一条指令分成五步后每条只需一个周期”，而是一个受资源和状态依赖约束的时间表。每个表格格子都应能回答：这条指令此刻占用什么资源、产生什么值、哪些后继可以消费、异常时哪些结果仍不可见。

\subsection{为什么常见五级划分只是一个可行模型}

经典 IF、ID、EX、MEM、WB 划分来自一组具体假设：取指、译码/寄存器读、算术逻辑、数据存储器访问和寄存器写回的组合延迟大致可分开，段间寄存器隔离各级状态。阶段数不是 ISA 事实；可以合并、拆分或复制阶段，只要保持语义等价和时序约束。段间寄存器带来额外建立/保持时间、填充与排空周期，也可能改变分支和数据依赖距离。

若各阶段延迟分别为 $t_i$、流水寄存器开销为 $t_r$，理想时钟周期近似为
\[
T_{clk}=\max_i(t_i)+t_r.
\]
一条指令的 latency 约为阶段数乘 $T_{clk}$，稳态 throughput 约为每周期一条，但实际总周期还要加入填充、排空、结构冲突、数据停顿、分支冲刷和 Cache/存储等待。

\subsection{结构冒险：资源冲突先于“相关性”}

结构冒险发生在同一周期需要同一个不可复制资源，例如单端口存储器同时取指和访问数据、单总线同时需要两个驱动、单写口寄存器堆同拍需要多次提交。解决路径只有几类：复制资源、增加端口、改变阶段调度、让一条指令 stall，或把资源使用改成可分离事务。题目问“为什么某周期必须停”，先画资源占用表，而不是先找寄存器依赖。

\subsection{数据冒险：依赖是静态关系，hazard 是时间冲突}

RAW、WAR、WAW 先描述程序中的读写依赖；只有当实际生产/消费时刻和资源调度造成冲突，才形成 hazard。简单顺序五级流水通常只暴露 RAW，因为读在前、写在后且顺序提交；乱序或更深流水可能出现 WAR/WAW，需要重命名、保留站或提交队列等额外机制。不能把所有“有依赖”都直接换成固定停顿数。

对每条相关边写四个时刻：
\[
\text{producer}\quad\text{produced}\quad\text{path-ready}\quad\text{consumer-need}.
\]
若 path-ready 早于或等于 need，可以 forwarding；若晚于 need，至少 stall 到可达。寄存器堆半拍读写、memory latency 和 forwarding 端点属于题设参数，不能从“经典五级”自动补齐。

\begin{center}\small
\begin{tabularx}{\linewidth}{L{2.0cm}L{3.0cm}L{3.0cm}Y}
\toprule
动作/问题 & 保存什么 & 作废什么 & 判据 \\
\midrule
forward & 已产生且可达的值 & 不改变年龄顺序 & producer 的 path-ready 不晚于 consumer-need \\
stall & PC/某些流水寄存器保持，插入 bubble & 不作废正确路径 & 等待资源或值变为可用 \\
flush & 年轻指令控制包清零/标无效 & 作废错误路径或异常后的年轻状态 & Next PC 已确定且错误路径不能提交 \\
\bottomrule
\end{tabularx}
\end{center}

\subsection{控制冒险与分支预测状态机}

分支冒险的代价由三个量生成：预测何时做、真实决策在哪一级完成、前端已经进入了多少错误路径指令。预测正确时只需继续或选择目标；预测错误时冲刷从预测点到决策点之间的年轻指令，并从正确 Next PC 重新取指。若题目给两位饱和计数器，其四状态通常按弱/强不跳与弱/强跳形成，正确结果向强方向移动，错误结果向另一方向移动；计数器状态不是“分支本身已执行”的架构状态。

固定“分支停两拍”只有在决策级、预测策略、前端深度和 flush 范围都被固定时才成立。分支目标生成、条件判断和异常优先级可能在不同阶段，必须从题设图重新数。

\subsection{流水线性能：吞吐量、延迟和总时间分开}

理想 $n$ 段流水执行 $m$ 条独立指令，常见总周期为 $n+m-1$；但这只是无停顿、每段单周期、无结构冲突的上界。实际可写成
\[
Cycles\approx IC+Fill+Drain+\sum Stall+\sum Flush+\sum Miss\ Penalty,
\]
再用 $CPU\ time=Cycles\times T_{clk}$ 或 $IC\times CPI\times T_{clk}$ 统一比较。吞吐量是稳态每单位时间完成多少条，latency 是单条指令从进入到结果/提交的时间；深流水可以改善频率和吞吐，却增加单条延迟、分支惩罚、依赖距离和寄存器开销。

超标量把同一周期可发射的独立指令数提高到大于一；多发射收益受真实 ILP、功能部件数量、寄存器端口和分支/Cache 等限制，不能把“理论发射宽度”直接当作实际 IPC。动态流水或乱序实现把 issue、execute、complete 和 commit 分开，结果可以先完成但仍须按年龄顺序提交，以保持精确异常。

\subsection{动态流水、硬件多线程与多核：并行对象必须先分类}

动态调度可以用保留站/评分表追踪操作数就绪和功能部件可用性，用重命名消除 WAR/WAW 的假依赖，再由 ROB 或等价结构按序提交。硬件多线程是在一个物理核心内保存多个线程上下文，交错发射以隐藏 cache miss 或长延迟，不等于增加 ALU 数量；多核则增加物理执行核心，进入 MIMD/共享内存层面的资源与一致性问题。

Flynn 分类只回答指令流/数据流的数量：SISD、SIMD、MISD、MIMD。SIMD/向量指令对多个数据元素使用同一控制，超标量是同一指令流的多条独立指令同拍发射，多核/SMP 是多个执行上下文或核心。题目出现“同时”时，先问新增的是数据元素、发射槽、线程上下文还是物理核心。

\subsection{异常和中断穿过流水线时的精确性}

同步异常由当前指令在某阶段产生，外中断由设备或外部事件提出；处理器可以在定义好的边界暂停前端、保存必要状态、选择入口并清空不能提交的年轻指令。故障类异常常允许修复后重新执行触发指令；陷阱、终止和外中断的返回点由 ISA 规定。关键不在“异常发生在哪一级”本身，而在：老指令哪些写回已提交、异常指令哪些副作用尚未提交、年轻指令哪些必须 flush、恢复后从哪一个 PC 重试。

\begin{methodbox}{异常题的三条线}
同时画 value dependency、ready/need timing 和 exception/commit 三条线。只画流水级名称，无法判断某个寄存器写回是否已经可见，也无法判断异常处理后是重试当前指令还是从下一条继续。
\end{methodbox}

\subsection{模拟器与时空图的验证用法}

流水线时空图或模拟器应按“预测—操作—回读”使用：
\begin{enumerate}
  \item 先不打开处理机制，记录无 forwarding/stall 时第一次冲突出现在哪一格；
  \item 只打开 forwarding，观察 ready 端点改变而不是直接删除依赖；
  \item 再打开 stall，记录哪一个 PC/流水寄存器保持、哪一个槽变成 bubble；
  \item 最后注入分支错误、异常或长延迟存储，核对 flush 范围、恢复 PC 和提交顺序。
\end{enumerate}
模拟器不是答案生成器；它的输出只有在能被 `producer/ready/need/commit` 四时刻解释时才算模型证据。

\section{扩充后的陌生流水线复原}

遇到新阶段、新预测规则或新发射宽度时，按以下顺序：
\begin{enumerate}
  \item 抄出每阶段资源、延迟、段间寄存器、读写时序和存储响应参数；
  \item 画无冲突时空图，确定填充、排空、结果可用和提交点；
  \item 对每条依赖标 producer、consumer、produced、path-ready、need、commit；
  \item 独立叠加结构冲突、RAW/WAR/WAW、分支预测和异常优先级；
  \item 对 stall 写“保持谁/插入什么”，对 flush 写“作废谁/从哪重取”；
  \item 统计完整周期并用 $IC\times CPI\times T_{clk}$ 检查数量级；
  \item 若问题进入多核共享 Cache、锁或操作系统调度，停止在本册并转到对应专题。
\end{enumerate}

\begin{methodbox}{CO-04 扩充停止条件}
读者能从陌生的阶段划分和预测规则重建时空图，解释一次 forwarding/stall/flush 的必要性，区分单条 latency 与稳态 throughput，并指出何时必须按序 commit 时，流水线机制层完成；具体缓存一致性、线程调度和编译器指令选择不在本册扩张。
\end{methodbox}

\section{补充细节：overview 增量：把性能指标放回 Cost 坐标}
性能题先区分两个目标：\kw{响应时间}是单个任务从提交到完成的墙钟时间，\kw{吞吐率}是单位时间完成的工作量。多核、批处理或深流水可能提高吞吐，却不必然缩短单个任务延迟；调度和 I/O 等待也会影响响应时间，但不应混入只问用户 CPU time 的公式。

CPU 执行时间的统一表达为
\[
T_{CPU}=IC\times CPI\times T_{clk}
       =\frac{IC\times CPI}{f},
\qquad
CPI=\frac{\text{总时钟周期}}{IC}.
\]
这里 $IC$ 受编译器、ISA 和程序语义共同影响，$CPI$ 受微架构、指令混合、Cache/TLB miss、分支和流水停顿影响，$T_{clk}$ 受关键路径、流水段划分、寄存器开销和实现工艺影响。三者不是可独立任意优化的旋钮：减少 $IC$ 的复杂指令可能增加 $CPI$，加深流水可能缩短 $T_{clk}$ 却增加 flush、旁路和寄存器成本。

\begin{center}\small
\begin{tabularx}{\linewidth}{L{2.4cm}L{3.0cm}L{3.1cm}Y}
\toprule
指标/量 & 衡量对象 & 常见公式或口径 & 不能直接推出\\\midrule
吞吐率 & 单位时间完成多少工作 & $\text{work}/\text{time}$ & 不等于单任务响应时间\\
响应时间 & 一个任务的完整墙钟耗时 & 提交至完成，含排队/等待 & 不等于用户 CPU time\\
IPS/MIPS & 每秒执行的指令数/百万指令数 & $IPS=f/CPI$ & 不同 ISA 的一条指令工作量不可直接等价\\
FLOPS & 每秒浮点操作数 & 由浮点操作计数定义 & 不能代表整数、I/O 或分支负载\\
CPU time & CPU 执行程序代码的时间 & $IC\times CPI\times T_{clk}$ & 不含墙钟 I/O 等待和其他任务占用\\
\bottomrule
\end{tabularx}
\end{center}

\subsection{Amdahl：局部加速的整体上限}
若原程序中比例 $p$ 的部分获得加速比 $s$，其余部分不变，则整体加速比为
\[
S_{overall}=\frac{1}{(1-p)+p/s}.
\]
当 $s\to\infty$ 时，$S_{overall}\le 1/(1-p)$；因此先判断未优化部分是否已构成上限，再决定是否继续堆硬件。若题目给出优化前后具体时间，先把“可优化部分”和“不可优化部分”分开，不要把局部 $s$ 当成整体加速。流水线、Cache、DMA、SIMD 和多核都只能加速它们实际覆盖的那部分工作。

\subsection{数量级与单位的防错}
主频/时钟周期使用十进制频率单位：$1\,\mathrm{GHz}=10^9\,\mathrm{Hz}$；地址空间和容量的幂次常按二进制解释：$1\,\mathrm{KiB}=2^{10}\,\mathrm{B}$。MIPS 的 “M” 是每秒百万条指令，FLOPS 的 “F” 是每秒浮点操作；题面写 MB、MiB、MHz 或 MT/s 时必须先确认它描述的是容量、基础时钟还是有效传输次数。

基准程序是负载集合而非单一程序。不同程序可能给出相反的机器排序，综合评价可用算术平均、几何平均或按使用频率加权；针对某个 benchmark 的特化优化也可能使结果不能代表一般负载。因而“主频更高”“MIPS 更大”或“流水级更多”都只能作为局部证据，最终仍回到统一的完整工作量和时间口径。

\begin{methodbox}{性能题的停止条件}
先声明问的是吞吐、响应时间、用户 CPU time 还是总线有效带宽；再列 $IC$、指令混合/CPI、$T_{clk}$ 和等待/重叠假设；最后用 Amdahl 上限或量纲检查结果。若只得到局部加速比而没有说明覆盖比例，性能结论尚未闭合。
\end{methodbox}

\subsection{并行扩展的细节边界：线程、核心与共享内存}
硬件多线程、超标量和多核都能让“更多工作同时推进”，但增加的对象不同：细粒度多线程可按周期轮换线程，粗粒度多线程通常在长延迟事件时切换，SMT 在同一周期从多个线程填充发射槽；它们共享同一物理核心的执行资源，不等于增加 ALU 数量。多核则增加物理执行核心，属于 MIMD 资源扩展。

共享内存多处理器还要区分：\kw{UMA} 中各核心访问统一内存的延迟近似一致；\kw{NUMA} 中延迟随目标内存节点而变，调度和数据放置会进入性能成本。SMP 的核心、Cache 和内存通过互连共享数据时，必须有一致性/同步协议保证不同核心对同一地址的可见性顺序；这不是单核 forwarding，也不能从“多核”自动推出某一种 MESI 状态机。题目只问 408 主干时，保留“共享副本需要一致性接口”这一边界，具体协议、锁和操作系统调度另行展开。

\begin{center}\small
\begin{tabularx}{\linewidth}{L{2.4cm}L{3.0cm}L{3.0cm}Y}
\toprule
对象 & 增加的并行维度 & 共享什么 & 主要 Cost/边界\\\midrule
SIMD/向量 & 一条指令处理多个数据元素 & 控制流与向量执行资源 & 数据相关/分支会降低利用率\\
超标量 & 同一线程同拍发射多条指令 & 核心内功能部件/寄存器 & 依赖、端口和提交宽度限制\\
硬件多线程 & 多个线程上下文交错/同时发射 & 一个物理核心资源 & 隐藏 miss，不自动增加物理算力\\
多核/SMP & 多个物理执行流 & 内存、互连和共享 Cache & 一致性、同步、UMA/NUMA 延迟\\
\bottomrule
\end{tabularx}
\end{center}

\end{document}
